\chapter{Represented Spaces and Type-Two Computation}
\label{sec:p04-represented-spaces}

Finite-input computation in \cref{def:partial-computable} computes a
natural number and halts. An arbitrary real number cannot be supplied
by a finite exact numeral. Type-Two Effectivity instead computes with
names that can provide indefinitely improving information.
The machine still manipulates symbols; the representation states
what object those symbols name.
The definitions below apply to arbitrary represented sets, not just
the reals \cite[pp.\ 4--18]{Weihrauch1995IntroductionRS}.

The source conventions remain identifiable. The held publisher
previews of Weihrauch's 2000 book cover pages 51--52 and 85--86,
not its complete naming-system or arithmetic chapters
\cite{Weihrauch2000AnalysisRS}.
The separately held July 1995 corrected second version and
Brattka--Hertling report supply the explicitly cited earlier
presentations \cite{Weihrauch1995IntroductionRS,BrattkaHertling1994RelationsRS}.
The Baire-stream presentation and full annex constructions below
are independently worded; they do not claim an unread 2000 proof route.

\section{Streams, names and productive machines}
\label{sec:p04-rs-names}

\begin{hbdefinition}{def:p04-rs-baire}{Baire streams and their finite observations}
Write \(\mathbb N=\{0,1,2,\ldots\}\) and
let \(\mathcal B=\mathbb N^{\mathbb N}\).
For a finite sequence \(s\in\mathbb N^{<\mathbb N}\), its cylinder is
\([s]=\{p\in\mathcal B:s\text{ is a prefix of }p\}\).
These cylinders form the basis of the product topology of discrete
\(\mathbb N\). Write \(p{\upharpoonright}k\) for the first \(k\)
entries of \(p\). Every domain \(D\subseteq\mathcal B\) below has
the subspace topology; no decidability or openness of \(D\) is presumed.
\end{hbdefinition}

A natural-valued stream is still a symbol input: encode it by
\(0^{p(0)}1\,0^{p(1)}1\cdots\), over a fixed finite alphabet.
Valid codes contain infinitely many delimiters.
Encoding and decoding each finite portion terminate, so a machine
can cache an input prefix and recover any already-read entry.
This is the Baire coding displayed in the held 2000
Definition 3.1.2(8), with the stream value rather than its binary
code used as the host name.
The finite-alphabet stream spaces in the 1995 presentation are
therefore not being identified with Baire space without a coding.

\begin{hbdefinition}{def:p04-rs-representation}{Representation and computable point}
A representation of a set \(X\) is a partial surjection
\(\delta_X:\mathcal B\rightharpoonup X\).
Its valid-name domain \(E_X=\dom(\delta_X)\) need not be computable.
If \(p\in E_X\) and \(\delta_X(p)=x\), then \(p\) is a name of \(x\).
Different names may denote the same object; invalid streams denote
no object. A represented space is \((X,\delta_X)\).
A point \(x\) is computable for this representation if it has at
least one computable name, meaning a finite program that produces
its entries in order indefinitely.
\end{hbdefinition}

In particular, being computable is not a property of the abstract
set \(X\) without a representation. Finite notations of countable
sets and infinite representations have distinct name types in the
original sources. A natural number can here be represented by
\(\delta_{\mathbb N}(p)=p(0)\), with irrelevant tail entries;
this is a convenient infinite-name counterpart of a finite notation,
not a claim that a halting finite output is an infinite computation.

\begin{hbdefinition}{def:p04-rs-machine}{Productivity and the machine domain}
A Type-Two machine has read-only input streams, finitely many work
tapes and an append-only output stream of finite-alphabet symbols.
It can perform arbitrarily long finite internal computations between
output symbols, but cannot erase an already issued symbol.
For Baire output it must complete the delimiter code of every entry.
Its productive domain consists of exactly those inputs on which
infinitely many complete output entries are produced.
On that domain its output defines a partial stream operator \(G\).
A finite output prefix or an endless computation that eventually
stops producing entries is not a completed Baire output.
No uniform time or input-lookahead bound is built into this definition.
\end{hbdefinition}

The source's finite-output variant halts with a finite word,
whereas its infinite-output variant computes forever and writes
the stream \cite[Definitions 2.1--2.3]{Weihrauch1995IntroductionRS}.
Neither a physical completion of an infinite run nor the
generalised Church--Turing thesis is a theorem asserted here.

\begin{hbdefinition}{def:p04-rs-realizer}{Realizer of a partial single-valued map}
For \(f:X\rightharpoonup Y\), put
\[
 D_f=\{p\in E_X:\delta_X(p)\in\dom(f)\}.
\]
A partial operator \(G:\mathcal B\rightharpoonup\mathcal B\)
realizes \(f\) if \(D_f\subseteq\dom(G)\), \(G(p)\in E_Y\)
for all \(p\in D_f\), and
\[
 \delta_Y(G(p))=f(\delta_X(p))\qquad(p\in D_f).
\]
The map is computable if a Type-Two machine's stream operator is
a realizer; it is continuously realizable if a continuous partial
stream operator is a realizer. All names of each valid input must
work, not just a chosen computable name.
No condition is imposed on \(G\) at invalid names or at names of
points outside \(\dom(f)\). Computability under this promise
convention does not decide whether an input satisfies the promise.
\end{hbdefinition}

For a single-valued \(f\), two names of the same input must yield
names of the same output object. This is decoded extensionality:
the two output streams need not be identical.
For a multivalued problem the outputs can denote different allowed
answers; that separate extension is not silently used to justify
a single-valued operation \cite[pp.\ 3--4]{BrattkaHertling1994RelationsRS}.

\begin{hbtheorem}{thm:p04-rs-finite-use}{proved=annex:p04-rs-finite-use}{Finite use, continuity and composition}
Every Type-Two machine operator is continuous on its productive
domain. Computable represented maps are consequently continuously
realizable. Compositions of computable represented maps are
computable, with their usual partial-map domain; the composed
machine may also be productive on unpromised inputs.
The same composition conclusion holds for continuous realizers.
\end{hbtheorem}

\section{Changing representations and topological admissibility}
\label{sec:p04-rs-admissibility}

\begin{hbdefinition}{def:p04-rs-reduction}{Translation and equivalence}
For two representations of the same set, write
\(\delta\leq_{\mathrm c}\gamma\) when there is a computable
translator \(T\) defined on all \(\dom(\delta)\) with
\(\gamma(T(p))=\delta(p)\).
Write \(\delta\leq_{\mathrm t}\gamma\) for a continuous translator.
Two representations are computably equivalent when each is
computably translatable into the other; continuous equivalence
is defined similarly. A translation cannot add information not
already computationally available from its input name.
\end{hbdefinition}

\begin{hbtheorem}{thm:p04-rs-invariance}{proved=annex:p04-rs-invariance}{Representation invariance with both directions}
Replacing input and output representations by computably
equivalent ones preserves computable points and computable
partial maps. Replacing them by continuously equivalent ones
preserves continuous realizability.
Continuous equivalence alone need not preserve computability.
\end{hbtheorem}

\begin{hbdefinition}{def:p04-rs-admissible}{Topological admissibility}
For a topological space \((X,\tau)\), a representation
\(\delta:E\to X\) is admissible here if it is continuous and has
this lifting property: for every subspace \(D\subseteq\mathcal B\)
and every continuous map \(h:D\to X\), there is a continuous
\(H:D\to E\) with \(\delta\circ H=h\).
Thus it accepts a continuous translation of every continuous
name-based presentation, including ones that cover only part of
\(X\). This is a topological universal property, not an assertion
that the translator has a finite program.
A quotient representation is a continuous surjection for which
\(U\subseteq X\) is open exactly when \(\delta^{-1}(U)\)
is open in \(E\).
\end{hbdefinition}

For the countably based \(T_0\) spaces developed next, this lifting
convention is equivalent to continuous equivalence with the
standard neighbourhood representation.
It is therefore compatible with the standard-representation
formulation in the held Brattka--Hertling introduction.
It is not being advertised as a full theorem about arbitrary
nonsequential spaces, or as the exact unread wording of
Weihrauch's 2000 admissibility chapter.

\begin{hbdefinition}{def:p04-rs-neighbourhood}{Names from a countable basis}
Let \((U_i)_{i\in\mathbb N}\) be a numbered countable basis of
a \(T_0\) space \(X\): every open set is a union of these basic
sets, and distinct points differ in membership in some open set.
Define
\[
 \delta_{\mathrm{nb}}(p)=x
 \quad\text{iff}\quad
 \{i:\exists n\ p(n)=i+1\}=\{i:x\in U_i\}.
\]
Zero entries are padding. Repetition and any order are allowed,
but every basic neighbourhood of \(x\), and no other one, must
occur. The \(T_0\) condition ensures uniqueness. Naming every
point is a set-theoretic property; its neighbourhood-index set
need not be computably enumerable.
\end{hbdefinition}

\begin{hbtheorem}{thm:p04-rs-standard}{proved=annex:p04-rs-standard}{Standard representations for general countably based spaces}
The neighbourhood representation is a continuous, admissible
quotient representation. A continuous representation of the same
space is admissible in the stated lifting sense if and only if it
is continuously equivalent to \(\delta_{\mathrm{nb}}\).
The quotient conclusion remains true when names are restricted
to any subspace \(A\subseteq X\) and the target has the subspace
topology.
\end{hbtheorem}

\begin{hbtheorem}{thm:p04-rs-topological}{proved=annex:p04-rs-topological}{Topological continuity versus continuous realization}
Suppose \(\delta_X:E_X\to X\) is a quotient representation and
\(\delta_Y:E_Y\to Y\) is continuous and admissible. For a total
single-valued \(f:X\to Y\), topological continuity is equivalent
to continuous realizability. For a partial \(f\), apply the same
assertion to its domain with its subspace topology, provided the
restricted input representation is quotient.
That last hypothesis holds for the neighbourhood representations
just constructed; it is not assumed for every quotient map.
\end{hbtheorem}

Admissibility supplies the continuous lifting direction.
The quotient condition supplies the descent from names to objects.
Computability further requires an effective translator or realizer;
neither an arbitrary numbering of a basis nor existence of a
continuous lift supplies that effectivity.

\section{Fast Cauchy names and their source bridges}
\label{sec:p04-rs-cauchy}

\begin{hbdefinition}{def:p04-rs-rational-code}{Effective rational data}
Fix the numerical pairing and its computable projections from
\cref{sec:turing-machines-01}. Define
\(\zeta(2a)=a\), \(\zeta(2a+1)=-a-1\), and the total rational
decoding
\[
 \nu_{\mathbb Q}(\operatorname{pair}_{\mathbb N}(a,b))
 =\frac{\zeta(a)}{b+1}.
\]
This is a surjection onto \(\mathbb Q\).
Finite signed fractions can be encoded, decoded and manipulated
by exact integer arithmetic; reduction by the Euclidean algorithm
gives a computable choice of a code for every rational.
Thus rational equality, order and the finite arithmetic operations
used below are effective. Distinct codes of one rational are allowed.
\end{hbdefinition}

\begin{hbdefinition}{def:p04-rs-cauchy}{Pointwise-error Cauchy representation}
For \(p\in\mathcal B\), put \(q_n=\nu_{\mathbb Q}(p(n))\).
Define
\[
 \delta_{\mathrm C}(p)=x
 \quad\text{iff}\quad
 x\in\mathbb R\ \text{and}\
 \lvert q_n-x\rvert\leq2^{-n}\quad\text{for every }n.
\]
Such an \(x\), if it exists, is unique. Rational density supplies
a name of every real. An arbitrary convergent sequence without
this supplied accuracy rate is not by itself a valid name.
No algorithm for recognizing the valid-name domain is assumed.
\end{hbdefinition}

Three related presentations must be distinguished.
First, let \((I_i)\) effectively list all open rational intervals.
Their enumeration representation \(\delta_{\mathrm I}\) names \(x\)
by listing exactly the indices \(i\) with \(x\in I_i\), using
positive index-plus-one entries and zero padding as before.
This is the one-dimensional neighbourhood-name construction.
The held 2000 page 86 explicitly describes names listing all
open rational intervals containing a point.

Second, the 1995 Definition 3.1 uses a pairwise Cauchy condition:
the decoded sequence \((u_n)\) names its limit \(x\) when
\[
 \lvert u_j-u_n\rvert<2^{-n}\qquad(j>n).
\]
Call this source-clause representation \(\rho_{\mathrm{pair}}\)
after changing only finite fraction syntax to rational codes.
Third, its interval representation \(\rho_{\mathrm{bounds}}\)
uses pairs of rational bounds \((\ell_n,v_n)\) satisfying
\(\sup_n\ell_n=\inf_n v_n=x\).
It does not require listing every open neighbourhood, nor does
its displayed definition require monotonically ordered bounds
\cite[pp.\ 19--20]{Weihrauch1995IntroductionRS}.

\begin{hbtheorem}{thm:p04-rs-cauchy-bridge}{proved=annex:p04-rs-cauchy-bridge}{Cauchy, neighbourhood and bound translations}
The four representations
\(\delta_{\mathrm C},\delta_{\mathrm I},
\rho_{\mathrm{pair}},\rho_{\mathrm{bounds}}\)
are computably equivalent, by translators valid on every name.
The pointwise-error representation \(\delta_{\mathrm C}\) is
continuous, admissible and quotient for the usual real topology;
its restriction to any represented real subspace is quotient
for that subspace topology.
These conclusions concern computability, not polynomial overhead.
\end{hbtheorem}

Together with \cref{thm:p04-rs-topological}, the quotient clause
implies that every Cauchy-computable partial real map is
topologically continuous on its domain with the subspace topology.
This does not make that domain recognizable or every continuous
real map computable.

The source's strict pairwise inequality is not the host's
non-strict error-to-limit inequality. A source name directly
satisfies the latter, but a host name needs a precision shift
to satisfy the former. The annex proves that shift and both
interval translations. Consequently an arithmetic algorithm
cannot be declared correct just because its outputs converge:
they must meet the actual output-name clause.

\section{Arithmetic algorithms and the nonzero promise}
\label{sec:p04-rs-arithmetic}

\begin{hbtheorem}{thm:p04-rs-arithmetic}{proved=annex:p04-rs-arithmetic}{Fast-name arithmetic}
With \(\delta_{\mathrm C}\) on real inputs and outputs and the
interleaved representation on pairs, addition, subtraction,
negation, absolute value, minimum, maximum and multiplication
are computable. On decoded names \(q_n\) of \(x\) and \(r_n\)
of \(y\), addition can output \(q_{n+1}+r_{n+1}\);
negation and absolute value can output \(-q_n\) and
\(\lvert q_n\rvert\).
For multiplication obtain \(B_x=\lvert q_0\rvert+1\) and
\(B_y=\lvert r_0\rvert+1\), and choose an integer \(m\) with
\[
 (B_x+B_y+1)2^{-m}\leq2^{-n}.
\]
Then \(q_mr_m\) is a valid \(n\)-th output approximation.
All choices and rational output codes are effective.
\end{hbtheorem}

The pair representation splits even and odd input entries into
the two names. Reading a requested entry needs only finitely
many stream symbols, so it supplies the product interface used
in that theorem.
The magnitude bound is obtained from the input itself.
For unbounded real arguments, multiplication is not justified
by a single precision offset independent of input magnitude.
The 1995 source gives algorithms for addition, multiplication,
maximum and reciprocal with its own pairwise output clause; the host theorem
also spells out the listed elementary extensions
\cite[Theorem 3.7(1), pp.\ 25--27]{Weihrauch1995IntroductionRS}.

\begin{hbtheorem}{thm:p04-rs-reciprocal}{proved=annex:p04-rs-reciprocal}{Reciprocal with two distinct accuracy clauses}
The partial operation \(x\mapsto1/x\), with domain
\(\mathbb R\setminus\{0\}\), is computable on fast Cauchy names.
For a host \(\delta_{\mathrm C}\)-name, search for \(k\) with
\[
 L=\lvert q_k\rvert-2^{-k}>0.
\]
After finding it, choose \(m\) for output position \(n\) with
\[
 2^{-m}\leq L/2,\qquad
 2^{1-m}/L^2\leq2^{-n},
\]
and output \(1/q_m\).
The search terminates on every name of a nonzero real.
On every valid name of zero it produces no output.

For the separate 1995 pairwise clause, search for \(N\) with
\(\lvert u_N\rvert>2\cdot2^{-N}\), and output
\(1/u_{2N+n}\).
This satisfies the strict pairwise output condition, not merely
convergence. Division \((x,y)\mapsto x/y\) is consequently
computable on the domain \(y\ne0\), by composition with
multiplication.
\end{hbtheorem}

This is one uniform program over all promised nonzero names:
the certificate and subsequent precision depend on the name.
It is not a program that decides whether an arbitrary real
equals zero, a total reciprocal at zero, or a uniform time bound
near zero. A known fixed lower bound on \(\lvert x\rvert\) can
remove the initial search; it is not silently supplied as input.

\begin{hbexample}{def:p04-rs-reciprocal-case}{A certificate is finite; its premise is not equality testing}
For the constant rational name \(q_n=1/8\), the first successful
host certificate occurs at \(k=4\), with \(L=1/16\).
Choosing \(m=n+9\) satisfies both displayed reciprocal bounds,
and the output is the rational \(8\) at every position.
For the constant source pairwise name \(u_n=1/8\), its different
certificate first succeeds at \(N=5\), and the source algorithm
uses \(u_{10+n}\). Both constructions name \(8\), under their
respective accuracy clauses.
The bounds are conservative because the algorithm cannot infer
from a finite prefix that all future approximants are exactly
equal. By contrast, a valid zero name always has
\(\lvert q_k\rvert\leq2^{-k}\), so no positive certificate appears.
Failure to find one after a finite wait is not evidence of equality.
\end{hbexample}

\section{Continuity is necessary, not computational completeness}
\label{sec:p04-rs-obstructions}

\begin{hbdefinition}{def:p04-rs-bit}{Discrete decision output}
Represent \(\{0,1\}\) by
\(\delta_2(p)=p(0)\) on streams whose first entry is zero or one.
The remaining entries are irrelevant. A total represented decision
must issue that first entry after finitely many steps on every
valid input name, then can pad its output indefinitely.
The induced topology is discrete.
\end{hbdefinition}

\begin{hbtheorem}{thm:p04-rs-zero}{proved=annex:p04-rs-zero}{Zero equality cannot be decided from a fast name}
The total characteristic function
\[
 Z(x)=
 \begin{cases}1,&x=0,\\0,&x\ne0,\end{cases}
\]
has no continuous \((\delta_{\mathrm C},\delta_2)\)-realizer
and hence no computable realizer.
There are nevertheless halting recognizers on valid Cauchy names
for \(x>0\), \(x<0\) and \(x\ne0\).
On the promised domain \(x\ne0\), the sign can be computed.
\end{hbtheorem}

The obstruction is a finite-information one. A finite zero-like
name prefix can still be extended to a name of a small nonzero
real. Increasing precision indefinitely therefore does not turn
all real predicates into total discrete decisions.
The semidecision searches use strict rational inequalities
\(q_k-2^{-k}>0\) or \(q_k+2^{-k}<0\), and wait forever
when their required strict separation never occurs
\cite[pp.\ 24--25]{Weihrauch1995IntroductionRS}.

\begin{hbtheorem}{thm:p04-rs-continuous-noncomputable}{proved=annex:p04-rs-continuous-noncomputable}{A continuous real map need not be computable}
There are real numbers with no computable Cauchy name.
For any such \(a\), the total constant map \(x\mapsto a\)
is topologically continuous and continuously realizable but
not computable in the displayed representations.
\end{hbtheorem}

Continuous dependence says that a finite required output
observation is fixed by a sufficiently informative finite input
observation at each valid input. It does not provide an effective
rule for selecting that observation or a finite program for
the output constants. In particular, pointwise existence of a
continuous lift is not a uniform computable assignment of lifts.

\begin{proofboundary}{def:p04-rs-extent}{Foundation and distinct source duties}
The present unit develops names, realizers, representation
translation, topological admissibility, real arithmetic and a
discontinuous decision obstruction with full local proofs.
The complete 2000 source presentation remains distinct from
its two held previews and the 1995 version.
Function-space operators such as differentiation, integration and
zero finding; Pour-El--Richards computability structures and
physical-equation constructions; represented complexity; and
ordinary/strong Weihrauch reductions and their degree operations
require their own hypotheses, representations and proofs.
None follows merely by calling a set a represented space.
\end{proofboundary}
